\chapter{Quantum Matter}

This chapter fully based on the content of lecture "Quantum Matter" taught by Professor Johannes Deiglmayr in Leipzig University in 2026.
\section{Review of atomic physics}
\subsection{Emission and Absorption of Electromagnetic Radiation by Atoms}
\paragraph{Indued and Spontaneous emission}
If an atom in the state $\ket k$ with energy $E_k$ is brought into an electromagnetic radiation field with spectral energy density $\omega_{\nu}(\nu)$ (field energy per unit volume and unit frequency interval $\Delta \nu = 1\mathrm{s}^{-1}$) it can absorb a photon $h \nu$, which brings the atom into a state with higher energy $E_{f}=E_{k}+ h \nu$.

The probability 




The energy of photon to excited the electron does not need to be exactly the same as the energy difference between the two levels, but there is a width of the absorption spectrum, which is called the natural linewidth. The natural linewidth is determined by the lifetime of the excited state, which is related to the uncertainty principle. The shorter the lifetime, the broader the linewidth. The natural linewidth can be calculated using the formula:
\begin{align}
\Gamma = \frac{\hbar}{\tau}
\end{align}
\section{Laser Cooling and Trapping}
\subsection{Deceleration of an Atomic Beam}

\section{Atomic Bose-Einstein Condensates and degenerate Fermi Gas}

\section{Optical Dipole Traps, optical lattice and application in quantum simulation}

\section{Precision measurements of fundamental constants/atomic clock}